\documentclass[11pt,a4paper]{article}

\usepackage[utf8]{inputenc}

\usepackage[T1]{fontenc}

\usepackage[polish]{babel}

\usepackage{lmodern}

\usepackage{geometry}

\geometry{a4paper,margin=27mm}

\usepackage{microtype}

\usepackage[hidelinks]{hyperref}

% Portfolio: MSC-AR-P11; source: prospective ARC-XR research scope.

\title{Metapoznanie w XR: informacja zwrotna o szacowanym obciążeniu}

\author{Mateusz Pomianek}

\date{Wersja 0.3.0, 21 listopada 2026}

\begin{document}

\maketitle

\begin{abstract}

Przedmiotem opracowania jest Metapoznanie w XR: informacja zwrotna o szacowanym obciążeniu. Czy informacja o szacowanym obciążeniu poprawia metapoznawczą kalibrację użytkownika? Protokół przewiduje dwie formy informacji zwrotnej: neutralny HUD i prezentację estymacji obciążenia, przy identycznych zadaniach poznawczych. Głównym punktem końcowym jest błąd kalibracji między pewnością deklarowaną a rzeczywistą poprawnością decyzji. Opis ma charakter wstępny i nie przedstawia niezweryfikowanych wyników.

\end{abstract}

\section{Wprowadzenie}

Rozwój adaptacyjnych systemów XR stwarza możliwość badania interakcji człowieka z informacją przestrzenną, ale samo zastosowanie AR, VR lub modeli sztucznej inteligencji nie przesądza o przewadze poznawczej czy ergonomicznej. W rozpatrywanym zagadnieniu, Metapoznanie w XR: informacja zwrotna o szacowanym obciążeniu, kluczowe jest rozdzielenie właściwości systemu od zmienności zachowania użytkownika. Porównanie należy zaprojektować tak, aby mierzyć jednoznaczne zjawisko, a nie tylko ogólną ocenę technologii.

\section{Cel pracy i pytanie badawcze}

Czy informacja o szacowanym obciążeniu poprawia metapoznawczą kalibrację użytkownika? Projekt wymaga krytycznego przeglądu aktualnego stanu badań oraz zdefiniowania nowości w stosunku do innych tematów portfolio XR.

\section{Zakres i odrębny wkład naukowy}
Wpływ informacji zwrotnej o oszacowanym obciążeniu na kalibrację subiektywnej pewności decyzji; skuteczność predykatora i przyczynowy efekt wyświetlenia są odrębne. Efekt sprawdzalnej informacji zwrotnej na kalibrację, nie diagnoza obciążenia albo sama trafność estymacji.
\section{Hipotezy badawcze}
\paragraph{H1.} Prawidłowo opisana informacja o obciążeniu zmniejsza błąd metapoznawczej kalibracji pewności względem neutralnego HUD.
\paragraph{H2.} Zysk nie występuje przy odwróconej lub nieinformatywnej informacji, co wskazuje efekt trafnej informacji, nie samej sugestii.
\paragraph{H3.} Skuteczność metapoznawcza nie pogarsza trafności zadania lub czasu wykonania poza ustalony margines.
\section{Projekt badania i plan eksperymentów}
\subsection{A. Operacjonalizacja i kwalifikacja danych}
Wpływ informacji zwrotnej o oszacowanym obciążeniu na kalibrację subiektywnej pewności decyzji; skuteczność predykatora i przyczynowy efekt wyświetlenia są odrębne. Jednostka wnioskowania: uczestnik × decyzja z deklarowaną pewnością i niezależnym wynikiem poprawności. W protokole zamrozić definicję i mianownik wyniku oraz źródło prawdy odniesienia, zgodność znaczników czasu, prawa do danych i zasady rekrutacji.
\subsection{B. Eksperyment porównawczy}
Losować neutralną, trafną i kontrolną nieinformatywną wskazówkę z jednakową widocznością; blokować wynik poprawności do niezależnej ewaluacji. Estymand pierwszorzędowy: różnica z góry określonego błędu kalibracji między informacyjnym i neutralnym HUD. Przed testem zdefiniować randomizację, próbę, punkt końcowy i warunki wykluczenia bez selekcji wyniku post hoc.
\subsection{C. Kontrole i analiza odporności}
Zabezpieczyć przed tym, aby badacz znał etykiety w chwili generowania wskazówek; wykazać wiarygodność estymatora na odrębnym materiale. Zmierzyć wpływ informacji na samoopis i wykrywać zgadywanie hipotezy przez uczestnika. W rejestrze uwzględnić braki, wstrzymania, opóźnienia i niepełne próby zamiast usuwać je z analizy.
\subsection{D. Walidacja na niezależnym materiale}
Weryfikować na niezależnych sesjach lub osobach, bez dostępu modelu do danych końcowych. Błąd kalibracji Brier lub wcześniej wybrany ECE z bootstrapem na uczestnikach; przed pilotażem zamrozić binning i margines. Dokumentować wyniki niesprzyjające hipotezie oraz rzeczywiste granice przeniesienia.
\section{Estymand, analiza statystyczna i decyzje}
różnica z góry określonego błędu kalibracji między informacyjnym i neutralnym HUD. Jednostka: uczestnik × decyzja z deklarowaną pewnością i niezależnym wynikiem poprawności. Błąd kalibracji Brier lub wcześniej wybrany ECE z bootstrapem na uczestnikach; przed pilotażem zamrozić binning i margines. H1 jest kandydatem hipotezy potwierdzającej, a H2/H3 pozostają eksploracyjne do formalnego planu wielokrotnych porównań. Liczebność uczestników, margines, próg i moc lub precyzja wymagają pilotażu oraz zamrożenia przed testem końcowym.
\section{Stan badań i granica nowości}
\begin{itemize}
\item Existing adaptive physiological feedback in AR means adapting an interface is not by itself a novel contribution. \cite{xrprior1}.
\item Explanation design in adaptive interfaces is prior art; response to veridical versus spurious explanations needs separate evaluation. \cite{xrprior2}.
\end{itemize}
Efekt sprawdzalnej informacji zwrotnej na kalibrację, nie diagnoza obciążenia albo sama trafność estymacji. Odrębność od wcześniejszych publikacji i innych tematów programu należy sprawdzić przed rejestracją badania.
\section{Warunki bezpieczeństwa, sprzętu i zgodności}
Jest to plan prospektywny, nie dowód wykonania eksperymentu lub zatwierdzonej etyki. Udział osób, rejestrowanie biometrii i interwencje wymagają uprzedniej procedury etycznej i świadomej zgody. Nie wnioskować o emocjach ani zdrowiu z samych sygnałów. Standardowy Quest 3 nie ma okulometrii HMD, a Tobii Eye Tracker 5 służy do stanowiska ekranowego; każde połączenie zweryfikować przed protokołem.
\section{Materiał i metody}

Protokół przewiduje dwie formy informacji zwrotnej: neutralny HUD i prezentację estymacji obciążenia, przy identycznych zadaniach poznawczych. Wymagane jest zdefiniowanie jednostki analizy, zasad doboru przypadków, wersji scenariusza i źródeł pomiaru. Rozkład warunków oraz reguły wykluczenia prób należy określić przed pomiarem. W badaniach z udziałem ludzi niezbędna jest uprzednia właściwa ocena etyczna i świadoma zgoda. W testach wyłącznie symulacyjnych należy oddzielić prawdę symulacyjną od obserwacji empirycznej.

\subsection{Wynik główny i kryteria jakości}

Błąd kalibracji między pewnością deklarowaną a rzeczywistą poprawnością decyzji. Pomiar powinien zachowywać jawne jednostki, znaczniki czasu i reguły oceny poprawności. Przy modelach predykcyjnych stosować rozłączne partycje oceny; przy eksperymentach interwencyjnych oddzielić skuteczność predykcji od przyczynowego wpływu interfejsu.

\section{Plan analizy}

Zestawić miary kalibracji i trafności przy uwzględnieniu uczestnika, trudności i kolejności prób. W wynikach należy raportować wielkość efektu lub błąd estymacji z przedziałem niepewności, liczbą obserwacji i odsetkiem odrzuconych prób. Analizy dodatkowe wymagają jasnego oznaczenia jako eksploracyjne. Nie wybierać progów decyzyjnych ani modelu na podstawie obserwacji ze zbioru testowego.

\section{Wyniki}

W obecnej wersji nie raportuje się wyników pomiarów, tabel statystycznych ani parametrów wyznaczonych z nieudokumentowanych danych. To kierunek badawczy, a nie zatwierdzony protokół badania ani zgłoszony artykuł.

\section{Dyskusja}

Interpretacja przyszłych wyników musi uwzględnić alternatywne mechanizmy i zależność wniosków od jakości danych, warunków zadania oraz ograniczeń sprzętu. Wartość potencjalnego rozwiązania zależy nie tylko od trafności modelu, lecz także od jego opóźnienia, stabilności i wpływu na zachowanie człowieka. Wyniki testów syntetycznych nie mogą być automatycznie uogólniane na populację uczestników.

\section{Ograniczenia badania}

Informacja zwrotna może być błędna i sama wpływać na decyzje; trafność estymatora wymaga niezależnej oceny. Granice nowości i trafności ekologicznej wymagają niezależnej weryfikacji. Żadna część szkicu nie dowodzi przeprowadzenia zbierania danych, uzyskania zgód etycznych ani gotowości manuskryptu do zgłoszenia.

\section{Wnioski}

Przedstawiono problem, kontrast i sposób jego planowanej oceny. Rozstrzygnięcie postawionego pytania wymaga przygotowania protokołu, zweryfikowanych źródeł oraz poprawnie udokumentowanych wyników. Bez tych materiałów nie można sformułować konkluzji empirycznych.

\begin{thebibliography}{9}
\bibitem{xrprior1} Designing adaptive feedback systems for managing cognitive load in augmented reality (2025), Design Science 11, e49. \url{https://doi.org/10.1017/dsj.2025.10040}.
\bibitem{xrprior2} EXAD: Framework for Designing Explanations in Adaptive User Interfaces (2025), IEEE Conference. \url{https://ieeexplore.ieee.org/abstract/document/11017085}.
\end{thebibliography}
\end{document}

